% !TEX program = pdflatex
% Example: Thesis using report class
% Compile with: pdflatex main.tex && bibtex main && pdflatex main.tex && pdflatex main.tex

\documentclass[12pt, a4paper]{report}

\usepackage{graphicx}
\usepackage{amsmath,amssymb}
\usepackage[margin=3cm]{geometry}
\usepackage{setspace}
\onehalfspacing
\usepackage{booktabs}
\usepackage{hyperref}

\begin{document}

% ---- Front matter ----
\title{Advanced Methods in Biomedical Signal Processing}
\author{Author Name}
\date{\today}
\maketitle

\begin{abstract}
This thesis presents novel methods for biomedical signal
processing with applications to brain-computer interfaces.
We develop new algorithms for feature extraction and
classification that improve performance across multiple
datasets.
\end{abstract}

\tableofcontents
\listoffigures
\listoftables

% ---- Main matter ----
\chapter{Introduction}
\section{Motivation}
Biomedical signal processing plays a crucial role in
modern healthcare. This thesis addresses the challenge
of extracting meaningful information from complex
biological signals.

\section{Contributions}
The main contributions are:
\begin{enumerate}
\item A novel feature extraction method.
\item An improved classification framework.
\item Extensive evaluation on public datasets.
\end{enumerate}

\chapter{Literature Review}
\section{Signal Processing Methods}
Traditional methods include Fourier analysis, wavelet
transforms, and empirical mode decomposition.

\chapter{Methodology}
\section{Proposed Framework}
Our framework combines time-frequency analysis with
deep learning architectures.

\chapter{Results}
\section{Experimental Setup}
We evaluated our method on three public datasets.

\section{Performance Analysis}
Results show significant improvement over baseline methods.

\chapter{Conclusion}
We presented a comprehensive framework for biomedical
signal processing. The proposed methods achieve
state-of-the-art performance.

% ---- Back matter ----
\bibliographystyle{plain}
\bibliography{references}

\appendix
\chapter{Mathematical Derivations}
\section{Proof of Theorem 1}
The proof follows from the properties of the wavelet
transform and the Cauchy-Schwarz inequality.

\end{document}
